\section{Unequal twists: an exact normal form and its finite reductions}
\label{tw:sec}

The incoming twisted-Stieltjes report asks for a normal form when two
frequency shifts are genuinely different
\cite[research question on unequal twists]{RepoTwisted}.
An equal-twist Gamma quotient cannot by itself answer that question.
We give three complementary answers: a convergent expansion in ordinary
Lerch families, an all-index convolution formula with one bounded
correction kernel, and a precise finite-reduction criterion for the
underlying scalar functions. The last criterion concerns functional
identities; it makes no inference about independence of sampled periods.

\subsection{Fourier conventions and the entire family}

Write $\T=\R/\Z$, $e_n(x)=e^{2\pi\ii nx}$, and
$\widehat U(n)=\langle U,e_{-n}\rangle$. Let $\theta,\eta$ be real
nonintegers, with $\theta\ne\eta$, and set
\begin{align}
 \lambda_{n,\theta}&=2\pi\ii(n+\theta),\\
 L_{n,\theta}&=\log(2\pi|n+\theta|)
                +\frac{\pi\ii}{2}\sgn(n+\theta).
 \label{tw:log}
\end{align}
Powers of $\lambda_{n,\theta}$ use this logarithm. The normalized
single-twist family and the mixed family are defined by
\begin{align}
 \widehat W_s^\theta(n)&=\lambda_{n,\theta}^{s-1},\label{tw:W}\\
 \widehat{\mathcal K_{\alpha,\beta}^{\theta,\eta}}(n)
   &=\lambda_{n,\theta}^{\alpha}\lambda_{n,\eta}^{\beta}.
 \label{tw:K}
\end{align}
Both are entire with values in periodic distributions, jointly in all
displayed spectral variables. Indeed their Fourier coefficients and every
fixed number of order derivatives have polynomial growth, uniformly on
compact spectral sets. Pairing with a smooth test function therefore gives
a normally convergent series. In particular,
\begin{equation}
 \mathcal K_{\alpha,\beta}^{\theta,\eta}
   =W_{1+\alpha}^{\theta}*W_{1+\beta}^{\eta}.
 \label{tw:Kconv}
\end{equation}
For $\Re(\alpha+\beta)<-1$, the Fourier series itself is absolutely
and uniformly convergent. The distributional definition supplies all
other spectral orders without assigning an ordinary value to a singular
endpoint.

For $0<\theta<1$, the standard Lerch family gives, on $0<x<1$,
\begin{equation}
 W_s^\theta(x)=\frac{e^{-2\pi\ii\theta x}}{\Gamma(1-s)}
       \Phi(e^{-2\pi\ii\theta},s,x).
 \label{tw:lerch}
\end{equation}
This identity is understood through analytic continuation when necessary.
It follows initially by unfolding the Fourier coefficient of
$\sum_{m\ge0}e^{-(\varepsilon+2\pi\ii\theta)(m+x)}(m+x)^{-s}$
and then letting $\varepsilon\downarrow0$. This is also the
normalization established in the incoming report. The Lerch function,
its relation $\Li_s(z)=z\Phi(z,s,1)$, and its analytic continuation
are classical \cite{DLMFLerch}. Other real twists follow by integer
modulation, $W_s^{\theta+k}=e_{-k}W_s^\theta$.

\subsection{A convergent Lerch expansion for arbitrary orders}

Put $d=\theta-\eta$ and $h=2\pi\ii d$. Choose a finite set
$F\subset\Z$ such that, for some $\rho<1$,
\begin{equation}
 \left|\frac{d}{n+\eta}\right|\le\rho\qquad(n\notin F).
 \label{tw:Fcondition}
\end{equation}
For example, all modes with $|n+\eta|\le2|d|$ may be put in $F$.
Let $\Pi_FU=\sum_{n\in F}\widehat U(n)e_n$.

\begin{theorem}[Normal convergence in single-twist Lerch coordinates]
\label{tw:tailtheorem}
For all $\alpha,\beta\in\C$,
\begin{align}
 \mathcal K_{\alpha,\beta}^{\theta,\eta}
 ={}&\sum_{n\in F}
        \lambda_{n,\theta}^{\alpha}\lambda_{n,\eta}^{\beta}e_n
       \notag\\
 &+\sum_{k=0}^{\infty}\binom\alpha k h^k
       (I-\Pi_F)W_{1+\alpha+\beta-k}^{\eta}.
 \label{tw:tailseries}
\end{align}
The series converges normally after pairing with any smooth test
function, locally uniformly in $(\alpha,\beta)$. Every fixed order
derivative in either spectral variable may be taken termwise. After
finitely many initial terms are removed, it converges normally in any
prescribed $C^q$ norm, locally on spectral parameter sets.
\end{theorem}

\begin{proof}
For $n\notin F$, condition \eqref{tw:Fcondition} implies that
$n+\theta$ and $n+\eta$ have the same sign. Therefore there is no
branch correction in
\[
 L_{n,\theta}=L_{n,\eta}
             +\log\left(1+\frac{d}{n+\eta}\right).
\]
The ordinary binomial theorem proves the equality of every Fourier
coefficient in \eqref{tw:tailseries}. On a compact set of $\alpha$,
$|\binom\alpha k|\le C(1+k)^M$ for suitable $C,M$. This elementary
bound follows either from the finite product or by Cauchy's formula
applied to $(1+z)^\alpha$ on a disc of radius between $\rho$ and
$1$; the latter also gives a geometric bound sufficient here.
Each summand is consequently bounded by a polynomial in $|n|$
times a summable geometric sequence in $k$. The test-function Fourier
coefficients absorb the polynomial. Cauchy's formula on a slightly
larger compact spectral set justifies all order derivatives.

For the $C^q$ assertion, take $k_0$ larger than
$q+1+\sup\Re(\alpha+\beta)$. Factor
$|\lambda_{n,\eta}|^{\Re(\alpha+\beta)-k_0}$ from the terms
with $k\ge k_0$. The remaining ratios are at most a fixed geometric
sequence, while the $q$ derivatives in $x$ still leave an absolutely
summable series over $n$. This proves the stated normal convergence.
\end{proof}

The finitely many removed modes are essential when the twists lie in
different integer intervals. Omitting them can change both the phase and
the value of the result. The theorem gives a convergently evaluable
normal form at every spectral order, including all Lerch order jets.

\subsection{One bounded kernel controls every mixed Stieltjes jet}

Let $P_\theta$ have Fourier coefficient
$\gamma+L_{n,\theta}$, and define
\begin{equation}
 B_\theta(u)=\Gamma(1-u)W_{1+u}^{\theta}
   =\delta_0+\sum_{m\ge0}\frac{(-1)^{m+1}}{m!}
                         U_m^\theta u^{m+1}.
 \label{tw:B}
\end{equation}
Thus $P_\theta=-U_0^\theta$. This agrees with the incoming
finite-part Lerch normalization. Introduce the difference
\begin{equation}
 A_{\theta,\eta}=P_\theta-P_\eta,
 \qquad\widehat A_{\theta,\eta}(n)=L_{n,\theta}-L_{n,\eta}.
 \label{tw:A}
\end{equation}

\begin{lemma}[The difference is an ordinary bounded kernel]
\label{tw:bounded}
If $\theta,\eta$ belong to the same open interval between consecutive
integers, then
\begin{equation}
 A_{\theta,\eta}(x)=2\pi\ii\int_\eta^\theta
       \frac{e^{-2\pi\ii t x}}{1-e^{-2\pi\ii t}}\dd t,
       \qquad 0<x<1.
 \label{tw:Aintegral}
\end{equation}
This function extends smoothly to each closed side of the endpoint,
is bounded, and has jump
\begin{equation}
 A_{\theta,\eta}(0+)-A_{\theta,\eta}(1-)
                    =2\pi\ii(\theta-\eta).
 \label{tw:Ajump}
\end{equation}
For arbitrary real noninteger endpoints, the same formula holds with
the path indented below the intervening integers. Equivalently, for
$\eta<\theta$, it is the Cauchy principal value of the real integral
plus $\pi\ii\sum_{k\in\Z\cap(\eta,\theta)}e_{-k}(x)$.
\end{lemma}

\begin{proof}
The Green kernel
\begin{equation}
 G_t(x)=\frac{e^{-2\pi\ii t x}}{1-e^{-2\pi\ii t}}
 \label{tw:G}
\end{equation}
has Fourier coefficient $1/(2\pi\ii(n+t))$. Integrating that
coefficient gives the logarithmic difference in \eqref{tw:A} if
no zero is crossed. Smoothness in $x$ up to each side follows by
differentiating the compact $t$-integral. Subtracting the endpoint
values gives \eqref{tw:Ajump}.

When $t$ crosses $-n$, continuation below the pole adds $+\pi\ii$
to the logarithmic difference, precisely the principal-logarithm phase
change from $-\pi/2$ to $+\pi/2$ in \eqref{tw:log}.
The residue of $2\pi\ii G_t(x)$ at an integer $k$ is $e_{-k}(x)$.
Subtracting those finitely many simple poles defines an ordinary bounded
function plus a finite Fourier polynomial. Its Fourier coefficients
are again exactly \eqref{tw:A}. Reversing orientation covers
$\theta<\eta$.
\end{proof}

For $0<\theta,\eta<1$, an alternative useful expression is
\begin{equation}
 A_{\theta,\eta}(x)
 =e^{-2\pi\ii\eta x}\Phi(e^{-2\pi\ii\eta},1,x)
  -e^{-2\pi\ii\theta x}\Phi(e^{-2\pi\ii\theta},1,x).
 \label{tw:ALerch}
\end{equation}
The two $1/x$ singularities cancel. If $\theta=p/q\in(0,1)$,
the corresponding term reduces to ordinary rational-shift digammas:
\begin{equation}
 e^{-2\pi\ii px/q}\Phi(e^{-2\pi\ii p/q},1,x)
 =-\frac{e^{-2\pi\ii px/q}}q
      \sum_{j=0}^{q-1}e^{-2\pi\ii pj/q}
                          \psi\left(\frac{x+j}{q}\right).
 \label{tw:rational}
\end{equation}
To prove it, split the Lerch defining sum into residue classes:
\[
 \Phi(e^{-2\pi\ii p/q},s,x)
   =q^{-s}\sum_{j=0}^{q-1}e^{-2\pi\ii pj/q}
                          \zeta\left(s,\frac{x+j}{q}\right).
\]
Continue this identity to $s=1$. The Hurwitz poles cancel because the finite character sum is
zero; $\gamma_0(a)=-\psi(a)$ gives the formula.

Write $\exp_*(uA)=\delta_0+\sum_{r\ge1}u^r A^{*r}/r!$.
Since $A$ is bounded and belongs to $L^1(\T)$, this exponential is
entire with values in $\delta_0+L^1(\T)$, by the usual convolution
norm estimate. Put
\begin{equation}
 C(u,v)=\frac{\Gamma(1-u)\Gamma(1-v)}{\Gamma(1-u-v)}.
 \label{tw:C}
\end{equation}

\begin{theorem}[All-index unequal-twist formula]
\label{tw:jets}
As analytic germs of periodic distributions,
\begin{equation}
 B_\theta(u)*B_\eta(v)
       =C(u,v)B_\eta(u+v)*\exp_*(uA_{\theta,\eta}).
 \label{tw:mixedB}
\end{equation}
Consequently every mixed finite-part product is the finite coefficient
formula
\begin{equation}
 U_m^\theta*U_n^\eta
  =(-1)^{m+n}m!n!\,[u^{m+1}v^{n+1}]
       C(u,v)B_\eta(u+v)*\exp_*(uA_{\theta,\eta}).
 \label{tw:alljets}
\end{equation}
Only powers $A^{*r}$ with $0\le r\le m+1$ occur. The $r=0$
part is the known equal-twist Gamma table; the other terms have Fourier
coefficients tending to zero and introduce no new point-supported
counterterm.
\end{theorem}

\begin{proof}
At a Fourier mode, the right side of \eqref{tw:mixedB} is
\[
 \Gamma(1-u)\Gamma(1-v)
      e^{(u+v)L_{n,\eta}}
      e^{u(L_{n,\theta}-L_{n,\eta})},
\]
which is exactly the left side. This proves the analytic identity by
normal convergence after testing. Taking the indicated coefficients
proves \eqref{tw:alljets}. For $r\ge1$, the Fourier coefficient of
$A^{*r}$ is $O(|n|^{-r})$, whereas a fixed spectral jet of
$B_\eta$ grows only as a fixed polynomial in $\log|n|$. Thus their
product tends to zero. A nonzero distribution supported at one point
is a finite sum of derivatives of its Dirac mass and has polynomial
Fourier coefficients; none can be added while preserving these
coefficients. This statement refers to the displayed canonical
decomposition, not to an arbitrary splitting of a distribution into
summands.
\end{proof}

The first symmetric identity is particularly compact:
\begin{equation}
 \boxed{\quad
 U_0^\theta*U_0^\eta
   =U_1^\theta+U_1^\eta-\zeta(2)\delta_0
                     -\frac12 A_{\theta,\eta}^{*2}.
 \quad}
 \label{tw:firstmixed}
\end{equation}
Indeed $2U_1^\theta=P_\theta^{*2}+\zeta(2)\delta_0$,
and polarization gives the result. Here $A^{*2}$ is an ordinary
continuous function, since its Fourier coefficients are $O(n^{-2})$.
The new term is therefore an explicitly evaluable ordinary convolution
of bounded Lerch differences, while the contact coefficient remains
fixed. An equal-twist formula with the correction omitted is false:
at mode $n$ it misses $\tfrac12(L_{n,\theta}-L_{n,\eta})^2$.

\subsection{Finite reductions and their precise boundary}

\begin{proposition}[Integer resolvent products]
\label{tw:partial}
For positive integers $p,q$, with $h=2\pi\ii(\theta-\eta)$,
\begin{align}
 \mathcal K_{-p,-q}^{\theta,\eta}
 ={}&\sum_{j=1}^p
  \frac{(-1)^q\binom{p+q-j-1}{p-j}}{h^{p+q-j}}
                  W_{1-j}^\theta\notag\\
 &+\sum_{j=1}^q
  \frac{(-1)^{q-j}\binom{p+q-j-1}{q-j}}{h^{p+q-j}}
                  W_{1-j}^\eta.
 \label{tw:partialformula}
\end{align}
If one of $\alpha,\beta$ is a nonnegative integer, there is likewise
a finite binomial reduction at the other twist. All integer pairs
therefore reduce to finitely many single-twist families. In particular,
\begin{equation}
 G_\theta*G_\eta=\frac{G_\eta-G_\theta}{2\pi\ii(\theta-\eta)}.
 \label{tw:resolvent}
\end{equation}
\end{proposition}

\begin{proof}
Expand $(A-h)^{-q}$ at $A=0$ and $(B+h)^{-p}$ at $B=0$, with
$A-B=h$. The displayed coefficients are exactly the principal parts
of $1/(A^pB^q)$ at its two poles. Their difference from that rational
function is an entire rational function vanishing at infinity and is
therefore zero. Substitute $A=\lambda_{n,\theta}$ and
$B=\lambda_{n,\eta}$. The polynomial cases follow by the finite
binomial theorem. Equation \eqref{tw:resolvent} is $p=q=1$.
\end{proof}

For positive integer $j$, the terms in \eqref{tw:partialformula}
are elementary on the interval:
\begin{equation}
 W_{1-j}^\theta(x)
  =\frac{(-1)^{j-1}}{(j-1)!(2\pi\ii)^{j-1}}
      \partial_\theta^{j-1}
      \frac{e^{-2\pi\ii\theta x}}{1-e^{-2\pi\ii\theta}}.
 \label{tw:elementary}
\end{equation}
The combined right side has its usual confluent limit
$W_{1-p-q}^{\theta}$ as $\eta\to\theta$, although individual
partial-fraction coefficients diverge. For example,
\begin{align}
 G_{1/4}*G_{3/4}(x)
  &=\frac{\ii}{2\pi}
    \left((1+\ii)e^{-3\pi\ii x/2}
                   -(1-\ii)e^{-\pi\ii x/2}\right),\label{tw:quarters}\\
 G_{1/2}^{*2}(x)&=e^{-\pi\ii x}
                         \left(\frac{x}{2}-\frac14\right).
 \label{tw:half}
\end{align}

\begin{theorem}[Finite scalar reduction criterion]
\label{tw:classification}
Fix distinct complex numbers $a,b$ and coherent branches on a simply
connected domain avoiding $-a,-b$. Consider
$f(z)=(z+a)^\alpha(z+b)^\beta$. It is a finite linear combination
of single-centre powers and order jets
\[
 (z+a)^r\Log^k(z+a),\qquad (z+b)^t\Log^\ell(z+b),
 \qquad k,\ell\in\N\cup\{0\},
\]
with arbitrary complex exponents and constant coefficients, if and only if
\begin{equation}
 \alpha\in\Z_{\ge0},\qquad\text{or}\qquad
 \beta\in\Z_{\ge0},\qquad\text{or}\qquad
 \alpha,\beta\in\Z.
 \label{tw:criterion}
\end{equation}
\end{theorem}

\begin{proof}
Sufficiency follows from the binomial theorem and rational partial
fractions. For necessity, suppose first $\alpha\notin\Z$ and
$\beta\notin\Z_{\ge0}$. Near $z=-a$, write $w=z+a$ and
$c=b-a\ne0$. Then
\[
 f(z)=w^\alpha\sum_{j\ge0}\binom\beta j c^{\beta-j}w^j.
\]
All infinitely many coefficients are nonzero. Terms based at $-b$
are holomorphic near $-a$. Apply analytic continuation once around
$w=0$ and subtract the original expression. On the left this
multiplies the displayed infinite series by
$e^{2\pi\ii\alpha}-1\ne0$; on the proposed right it removes every
term holomorphic there and leaves only finitely many powers times
polynomials in $\Log w$. Let $E=w\,d/dw$. A nonzero polynomial $P(E)$ annihilates that
finite expression: for each occurring $w^r\Log^k w$, include the
factor $(E-r)^{k+1}$. Applied to the left, it gives a nonzero constant
multiple of
\[
 w^\alpha\sum_{j\ge0}P(\alpha+j)
                    \binom\beta j c^{\beta-j}w^j.
\]
Vanishing would force $P(\alpha+j)=0$ for every $j\ge0$, contrary
to $P$ being a nonzero polynomial. This is a contradiction.

Interchanging the two centres proves the same conclusion if
$\beta\notin\Z$ and $\alpha\notin\Z_{\ge0}$. Unless
\eqref{tw:criterion} holds, at least one of these two cases occurs.
\end{proof}

The qualification ``scalar'' matters. Equality on the discrete set of
Fourier modes is weaker than equality of the underlying holomorphic
functions; the theorem does not promote the former to the latter.
It explains exactly when the elementary functional mechanism behind
\eqref{tw:partialformula} can terminate. Arbitrary orders still have
the normally convergent expansion \eqref{tw:tailseries}.

\subsection{Argument derivatives and canonical primitives}

Let $D_\theta=d/dx+2\pi\ii\theta$. Fourier multiplication gives
the exact contiguous and differential laws
\begin{align}
 D_\theta\mathcal K_{\alpha,\beta}^{\theta,\eta}
   &=\mathcal K_{\alpha+1,\beta}^{\theta,\eta},&
 D_\eta\mathcal K_{\alpha,\beta}^{\theta,\eta}
   &=\mathcal K_{\alpha,\beta+1}^{\theta,\eta},\label{tw:diff}\\
 \mathcal K_{\alpha+1,\beta}^{\theta,\eta}
  -\mathcal K_{\alpha,\beta+1}^{\theta,\eta}
   &=h\mathcal K_{\alpha,\beta}^{\theta,\eta}.
 \label{tw:contiguous}
\end{align}
Since neither twist is integral, $D_\theta$ and $D_\eta$ are
invertible on periodic distributions. Their inverses lower the
corresponding spectral exponent. These statements remain valid after
any order derivatives in $\alpha,\beta$.

For an ordinary integrable $f$, the unique periodic covariant primitive
has the explicit form
\begin{equation}
 (D_\eta^{-1}f)(x)=e^{-2\pi\ii\eta x}
 \left\{\int_0^x e^{2\pi\ii\eta t}f(t)\dd t
   +\frac{e^{-2\pi\ii\eta}}{1-e^{-2\pi\ii\eta}}
          \int_0^1e^{2\pi\ii\eta t}f(t)\dd t\right\}.
 \label{tw:primitive}
\end{equation}
Differentiation proves $D_\eta U=f$, and the displayed constant is
forced by $U(1)=U(0)$. This supplies ordinary quadratures for every
bounded correction in \eqref{tw:alljets}, while the singular
single-twist part retains its canonical distributional primitive.

